\section{Local genus-one models}\label{sec:local-genus-one}
We recall the comparison of relative differentials from \cite{LLR04} and prove a congruence used in the genus-one case of the main theorem.
%We call regular genus-one models over discrete valuation rings \emph{surfaces}.

\subsection{Local models and relative differentials}\label{subsec:genus-one-setting}
Let $R$ be a discrete valuation ring with uniformizer $\pi$, residue field $k = R/\pi R$, and fraction field $K$.
Let $Y$ be a smooth projective curve of genus one over $K$ with $H^0(Y,\O_Y) = K$.
Let $J$ be its Jacobian.
Let $\mathcal Y$ and $\mathcal J$ be the {\em minimal regular models} of $Y$ and $J$ over $R$, respectively; see for example \cite[Section~9.3]{Liu02}.
We say that $J$ has \emph{good reduction} over $R$ if it extends to an elliptic scheme over $R$, equivalently if $\mathcal J$ is smooth over $R$.
There is a canonical cohomology isomorphism \cite[1.4]{LLR04}
\begin{equation}\label{eq:tau}
  \tau_Y\from H^1(Y,\O_Y)\xrightarrow{\sim}H^1(J,\O_J).
\end{equation}
This isomorphism is compatible with field extension and, when $Y(K)\neq\varnothing$, is induced by a choice of $K$-point identifying $J$ with $Y$ \cite[1.4 and Corollary~1.5]{LLR04}.
The comparison theorem \cite[Theorem~3.8]{LLR04} shows that $\tau_Y$ extends to an $R$-linear homomorphism
\[
  \tau_{\mathcal Y}\from H^1(\mathcal Y,\O_{\mathcal Y})\to H^1(\mathcal J,\O_{\mathcal J}).
\]

The target is a free $R$-module of rank one \cite[Corollary~5.8]{LLR04}.
Since $\tau_{\mathcal Y}$ becomes an isomorphism over $K$, its kernel is exactly the $R$-torsion submodule $H^1(\mathcal Y,\O_{\mathcal Y})_{\rm tors}$ of $H^1(\mathcal Y,\O_{\mathcal Y})$, and its cokernel is also torsion.
Both have finite $R$-length because the cohomology groups are finitely generated over $R$.

\begin{definition}\label{def:hodge}
We define
\[
  a(\mathcal Y) = \length_R\ker(\tau_{\mathcal Y}),\qquad b(\mathcal Y) = \length_R\coker(\tau_{\mathcal Y}).
\]
\end{definition}

By relative Serre duality, $H^0(\mathcal Y,\omega_{\mathcal Y/R})$ is the $R$-dual of $H^1(\mathcal Y,\O_{\mathcal Y})$, and similarly for $\mathcal J$; these modules are free of rank one.
Let $\eta_{\mathcal Y}$ and $\eta_J$ be generators of $H^0(\mathcal Y,\omega_{\mathcal Y/R})$ and $H^0(\mathcal J,\omega_{\mathcal J/R})$, respectively, and identify $H^0(Y,\omega)$ with $H^0(J,\omega)$ by the dual of $\tau_Y$.
Write $\mathcal Y_s$ for the closed fibre of $\mathcal Y$.

\begin{lemma}\label{lem:ab-properties}
\begin{enumerate}
\item $a(\mathcal Y) = 0$ if and only if $H^0(\mathcal Y_s,\O_{\mathcal Y_s}) = k$.
\item If $\mathcal Y$ is smooth over $R$, then $a(\mathcal Y) = b(\mathcal Y) = 0$.
\item $\mathcal Y\otimes_R\widehat R$ and $\mathcal J\otimes_R\widehat R$ are the minimal regular models over the completion $\widehat R$, and $a(\mathcal Y\otimes_R\widehat R) = a(\mathcal Y)$, $b(\mathcal Y\otimes_R\widehat R) = b(\mathcal Y)$.
\item If $k$ is perfect, or $J$ has split multiplicative reduction, then $a(\mathcal Y) = b(\mathcal Y)$.
%If $k$ is imperfect, the equality $a(\mathcal Y) = b(\mathcal Y)$ can fail \cite[Theorem~5.9 and Proposition~9.2]{LLR04}.
\item \textup{(Comparison of relative differentials)}
  The inclusion $H^0(\mathcal J,\omega_{\mathcal J/R})\embed H^0(\mathcal Y,\omega_{\mathcal Y/R})$ has cokernel of $R$-length $b(\mathcal Y)$.
  In particular,
  \begin{equation}\label{eq:hodge}
    \eta_J = \pi^{b(\mathcal Y)}\eta_{\mathcal Y}\quad\text{up to a unit of } R.
  \end{equation}
  This identity persists in $H^0(Y_{K'},\omega)$ after any flat base change $R\to R'$ of discrete valuation rings, where $K' = \Frac(R')$.
\end{enumerate}
\end{lemma}

\begin{proof}
(1) Since $\mathcal Y$ is integral and dominates $\Spec R$, multiplication by $\pi$ is injective on $\O_{\mathcal Y}$.
The ring $H^0(\mathcal Y,\O_{\mathcal Y})$ is a finite $R$-subalgebra of $H^0(Y,\O_Y) = K$, so it equals $R$ because $R$ is normal.
The exact sequence $0\to\O_{\mathcal Y}\xrightarrow{\pi}\O_{\mathcal Y}\to\O_{\mathcal Y_s}\to 0$ therefore gives
\[
  0\to k\to H^0(\mathcal Y_s,\O_{\mathcal Y_s})\to H^1(\mathcal Y,\O_{\mathcal Y})[\pi]\to 0.
\]
A finitely generated $R$-module is torsion free if and only if its $\pi$-torsion vanishes.

(2) After a faithfully flat unramified base change, the smooth model $\mathcal Y$ has a section and is therefore isomorphic to $\mathcal J$.
By flat base change and the description of $\tau_Y$ above, $\tau_{\mathcal Y}$ becomes an isomorphism; faithful flatness then gives $a(\mathcal Y) = b(\mathcal Y) = 0$.

(3) The scheme $\mathcal Y\otimes_R\widehat R$ is regular \cite[Tag~\stackstag{0BG6}]{Stacks} and has the same closed fibre as $\mathcal Y$, so it is the minimal regular model over $\widehat R$, and similarly for $\mathcal J$.
Flat base change identifies $\tau_{\mathcal Y\otimes_R\widehat R}$ with $\tau_{\mathcal Y}\otimes_R\widehat R$, whose kernel and cokernel have the same lengths over $\widehat R$ as those of $\tau_{\mathcal Y}$ over $R$.

(4) The perfect-residue-field case is \cite[Theorem~3.8]{LLR04}.
By (3), the split multiplicative case reduces to \cite[Remark~8.13]{LLR04} over $\widehat R$.

(5) The map induced by $\tau_{\mathcal Y}$ on $H^1(\mathcal Y,\O_{\mathcal Y})/{\rm tors}$ is injective, with cokernel of length $b(\mathcal Y)$.
Dualizing gives an inclusion of the Jacobian differential lattice $H^0(\mathcal J,\omega_{\mathcal J/R})$ into $H^0(\mathcal Y,\omega_{\mathcal Y/R})$ with the same cokernel length, which is \eqref{eq:hodge} for the chosen generators.
Since $\tau_Y$ is compatible with base change \cite[1.4]{LLR04}, the identity persists after flat base change.
\end{proof}

\subsection{Good Jacobian reduction}\label{subsec:good-jacobian-reduction}
Assume further that $R$ is excellent of equal characteristic $p>0$ and that $J$ has good reduction.
The curve $Y$ is a torsor under $J$.
By \cite[Proposition~8.1 and Remark~8.2]{LLR04}, this action extends to $\mathcal J\times_R\mathcal Y\to\mathcal Y$ and is transitive on the geometric points of the closed fibre.
In particular, the closed fibre $\mathcal Y_s$ is geometrically irreducible.
As an effective Cartier divisor, it has the form
\[
  \mathcal Y_s=rV,
\]
where $r\geq1$ is an integer and $V$ is integral.
By \cite[Proposition~8.1]{LLR04}, $V$ is smooth of genus one over $k_0=H^0(V,\O_V)$, and $k_0/k$ is finite and purely inseparable.

Let $\eta_{\mathcal Y}$ be a generator of $H^0(\mathcal Y,\omega_{\mathcal Y/R})$ as before.
Its restriction to the smooth generic curve has no zeros, so its zero divisor is $qV$ for a unique integer $q\geq0$.
If $q\geq r$, then $(1/\pi)\eta_{\mathcal Y}$ has divisor $(q-r)V\geq0$ and hence is a global regular section, contradicting $H^0(\mathcal Y,\omega_{\mathcal Y/R})=R\eta_{\mathcal Y}$.
Thus
\begin{equation}\label{eq:local-canonical-order}
  0\leq q<r,
  \qquad \div(\eta_{\mathcal Y})=qV.
\end{equation}

Via the identification of $H^0(Y,\omega)$ with $H^0(J,\omega)$ in Subsection~\ref{subsec:genus-one-setting} and Lemma~\ref{lem:ab-properties}(5), we regard the Jacobian differential $\eta_J$ as a global section of $\omega_{\mathcal Y/R}$.
By \eqref{eq:hodge}, $\eta_J = \pi^{b(\mathcal Y)}\eta_{\mathcal Y}$ up to a unit; since $\div(\pi) = rV$, \eqref{eq:local-canonical-order} gives $\div(\eta_J) = \chi V$, where
\[
  \chi = r\,b(\mathcal Y)+q.
\]
Since $0\leq q<r$, the integer $\chi$ determines both $b(\mathcal Y) = \lfloor\chi/r\rfloor$ and $q$.
The integer $q$ plays the same role as the coefficient $a$ in the canonical bundle formula \cite[Theorem~2]{BM77}.
When the residue field is algebraically closed, $\chi$ is computed by the following result of Raynaud, whose proof is written out in \cite{BT14}.

\begin{proposition}[{\cite[Proposition~3.17]{BT14}}]\label{prop:raynaud-chi}
Assume that $R$ is complete and $k$ is algebraically closed, and let $\nu$ be the order of $\O_{\mathcal Y}(V)|_V$ in $\Pic(V)$.
Then $r = \nu p^s$ for an integer $s\geq0$, and there are positive integers $k_0,\dots,k_{s-1}$ such that
\begin{equation}\label{eq:raynaud-chi}
  \chi = (r-1)+\sum_{i=0}^{s-1}k_i\,(r-\nu p^i).
\end{equation}
\end{proposition}

Here $s$ and the $k_i$ record how the order of $\O_{\mathcal Y}(-V)|_{nV}$ in $\Pic(nV)$ grows from $\nu$ to $r$ along the thickenings $nV\subset\mathcal Y$ defined by $\O_{\mathcal Y}(-nV)$ \cite[Lemma~3.12]{BT14}.
In the tame case $s = 0$ we get $\chi = r-1$, so $b(\mathcal Y) = 0$ and $q = r-1$; if $s\geq1$, then $\chi\geq r$ and $b(\mathcal Y)\geq1$.
Since $r-\nu p^i = \nu p^i(p^{s-i}-1)$ is divisible by $p-1$, \eqref{eq:raynaud-chi} gives $\chi\equiv r-1\pmod{p-1}$.
After passing to the completion of the strict henselization of $R$, this implies Theorem~\ref{thm:local-hodge-congruence} below when $k$ is perfect; the rest of the section treats imperfect $k$ by reduction to this case.

\begin{theorem}\label{thm:local-hodge-congruence}
In this setting,
\begin{equation}\label{eq:local-hodge-congruence}
  r\,b(\mathcal Y)+q\equiv r-1\pmod{p-1}.
\end{equation}
\end{theorem}

\begin{corollary}\label{lem:congruence}
Assume that $r = 1$, so that $\mathcal Y_s = V$ is reduced.
Then $b(\mathcal Y)$ is divisible by $p-1$.
Moreover, $\mathcal Y$ is smooth over $R$ if and only if $b(\mathcal Y) = 0$.
\end{corollary}

If $\mathcal Y_s$ is reduced and $\mathcal Y/R$ is not smooth, then the image of $\eta_J$ in $H^0(\mathcal Y,\omega_{\mathcal Y/R})$ vanishes to order at least $p-1$ along $\mathcal Y_s$ by \eqref{eq:hodge}.

\subsection{Normalization under extension of constants}
To prove Theorem~\ref{thm:local-hodge-congruence}, we will compare $\mathcal Y$ with a regular model over a discrete valuation ring with algebraically closed residue field.
Extending the coefficient field can destroy regularity, so we must normalize after base change.
We first treat finite extensions and compute the changes in $r$, $q$, and $b$; these relations will then allow us to construct a finite regular normalization after extension to $\bar k$.

\begin{lemma}\label{lem:local-normalization}
Let $\mathcal Y/R$ be as in Subsection~\ref{subsec:good-jacobian-reduction}, and assume moreover that $R$ is henselian with a separably closed coefficient field $k$.
Let $\ell/k$ be a finite purely inseparable extension and put $R_{\ell} = R\otimes_k \ell$.
\begin{enumerate}
\item The normalization $\mathcal Y_{\ell}$ of $\mathcal Y\times_R\Spec R_{\ell}$ is its minimal regular model.
\item Let $\ell'/\ell$ be a further finite purely inseparable extension.
Write $r_{\ell},V_{\ell},q_{\ell}$ for $r,V,q$ of Subsection~\ref{subsec:good-jacobian-reduction} applied to $\mathcal Y_{\ell}$, and similarly for $\ell'$.
There are nonnegative integers $e = e(\ell'/\ell)$, $N = N(\ell'/\ell)$ and $n = n(\ell'/\ell)$, with $(p-1)\mid N$, such that
\begin{align}
  r_{\ell'}&=p^e r_{\ell},\label{eq:local-multiplicity-change}\\
  n r_{\ell'}&=N+q_{\ell'}-p^e q_{\ell},\label{eq:local-canonical-change}\\
  b(\mathcal Y_{\ell})&=n+b(\mathcal Y_{\ell'}).\label{eq:local-hodge-change}
\end{align}
Here $p^e$ is the ramification index along $V_{\ell'}$ over $V_{\ell}$, and the divisorial conductor of $\mathcal Y_{\ell'}\to(\mathcal Y_{\ell})_{R_{\ell'}}$ is $NV_{\ell'}$.
\end{enumerate}
\end{lemma}

\begin{proof}
(1) Since $R_{\ell}$ is finite free over $R$ and $R_{\ell}/\pi R_{\ell}\cong\ell$, it is a one-dimensional Noetherian local ring with maximal ideal $\pi R_{\ell}$, hence a discrete valuation ring with uniformizer $\pi$.
The ring $R_{\ell}$ is excellent because it is finite over $R$, so normalization is finite.
Since the generic curve is geometrically integral, the base change is integral.
The composite $\mathcal Y_{\ell}\to\mathcal Y$ is finite and purely inseparable with normal target, hence a universal homeomorphism; so is the normalization map.

The action of $\mathcal J_{R_{\ell}}$ on $\mathcal Y\times_R\Spec R_{\ell}$ lifts to $\mathcal Y_{\ell}$.
The product $\mathcal J_{R_{\ell}}\times_{R_{\ell}}\mathcal Y_{\ell}$ is normal because it is smooth over the normal scheme $\mathcal Y_{\ell}$.
Hence the induced morphism to $\mathcal Y\times_R\Spec R_{\ell}$ factors through its normalization $\mathcal Y_{\ell}$.
The identity and associativity axioms hold on the generic fibre and hence everywhere.
The action on geometric points of the closed fibre is transitive because the normalization map is a universal homeomorphism.

Let $S$ be the singular locus of the normal excellent surface $\mathcal Y_{\ell}$; it is finite, lies in the closed fibre, and is invariant under translations by $\mathcal J(R_{\ell})$.
Every $\ell$-point of $\mathcal J_s$ lifts to an $R_{\ell}$-point because $R_{\ell}$ is henselian and $\mathcal J$ is smooth.
Since $\ell$ is separably closed, the group $G=\mathcal J_s(\ell)$ is infinite.

Fix a closed point $x\in(\mathcal Y_{\ell})_s$ and a geometric point $\bar x$ over it.
By transitivity, the geometric orbit map $g\mapsto g\cdot\bar x$ is a nonconstant morphism from an elliptic curve, so its fibres are finite.
Hence the subgroup $H=\{g\in G\mid g\cdot\bar x=\bar x\}$ is finite, and the orbit $G\cdot\bar x\cong G/H$ is infinite.
Since $\ell$ is separably closed, each closed point has a unique geometric point over it, so these orbit points lie over infinitely many distinct closed points.
The lifting statement above places all these closed points in the $\mathcal J(R_{\ell})$-orbit of $x$.
If $x$ were singular, its entire orbit would lie in the finite set $S$, a contradiction.
Thus $S$ is empty and $\mathcal Y_{\ell}$ is regular.
The reduced closed fibre is a homogeneous curve under an elliptic curve.
The argument of \cite[Proposition~8.1]{LLR04} shows that it is smooth of genus one over its purely inseparable constant field.
It contains no $(-1)$-curve, so $\mathcal Y_{\ell}$ is minimal.

(2) Put $\mathcal X=\mathcal Y_{\ell}\times_{R_{\ell}}\Spec R_{\ell'}$.
\[
  \begin{tikzcd}
    \mathcal Y_{\ell'} \arrow[r,"\nu"] & \mathcal X \arrow[r] \arrow[d] & \mathcal Y_{\ell} \arrow[d]\\
    & \Spec R_{\ell'} \arrow[r] & \Spec R_{\ell}.
  \end{tikzcd}
\]
The right square is Cartesian, and $\nu$ is the normalization map.
The composite $\mathcal Y_{\ell'}\to\mathcal Y_{\ell}$ is finite and purely inseparable, so its ramification index along $V_{\ell'}$ is $p^e$ for some integer $e\geq0$.
Since $\pi$ is a uniformizer of both $R_{\ell}$ and $R_{\ell'}$, the fibre multiplicities satisfy
\[
  r_{\ell'} = \operatorname{ord}_{V_{\ell'}}(\pi)
  = p^e\operatorname{ord}_{V_{\ell}}(\pi) = p^e r_{\ell},
\]
which proves \eqref{eq:local-multiplicity-change}.
The generic fibre is smooth, so the divisorial conductor is $NV_{\ell'}$ for an integer $N\geq0$.
The regular model $\mathcal Y_{\ell}/R_{\ell}$ is Gorenstein.
For every finite extension $\ell''/\ell$, the scheme $\mathcal Y_{\ell}\times_{\ell}\Spec\ell'' = \mathcal Y_{\ell}\times_{R_{\ell}}\Spec R_{\ell''}$ is integral because it is flat over the discrete valuation ring $R_{\ell''}$ with generic fibre a base change of the geometrically integral curve $Y$.
Since $\ell$ is separably closed, $\mathcal Y_{\ell}$ is therefore geometrically integral over $\ell$ \cite[Tags~\stackstag{035X} and~\stackstag{038K}]{Stacks}.
Applying Lemma~\ref{lem:canonical-base-change} with $(k,A,X,k') = (\ell,R_{\ell},\mathcal Y_{\ell},\ell')$ gives
\[
  \omega_{\mathcal Y_{\ell'}/R_{\ell'}}
  \cong\nu^*\omega_{\mathcal X/R_{\ell'}}\otimes\O_{\mathcal Y_{\ell'}}(-NV_{\ell'}).
\]
In particular, $(p-1)\mid N$.

Let $\eta_{\mathcal Y_{\ell}}$ and $\eta_{\mathcal Y_{\ell'}}$ be the generators of Subsection~\ref{subsec:good-jacobian-reduction} for $\mathcal Y_{\ell}$ and $\mathcal Y_{\ell'}$, and write $\eta_{\mathcal Y_{\ell}}=\pi^{-n}\eta_{\mathcal Y_{\ell'}}$ up to a unit.
By \eqref{eq:local-canonical-order} applied to $\mathcal Y_{\ell}$, we have $\div(\eta_{\mathcal Y_{\ell}}) = q_{\ell}V_{\ell}$.
Since $\omega_{\mathcal X/R_{\ell'}}$ is the pullback of $\omega_{\mathcal Y_{\ell}/R_{\ell}}$ and $V_{\ell}$ pulls back to $p^eV_{\ell'}$ on $\mathcal Y_{\ell'}$, the pullback of $\eta_{\mathcal Y_{\ell}}$ is a section of $\nu^*\omega_{\mathcal X/R_{\ell'}}$ with divisor $p^e q_{\ell}V_{\ell'}$.
Via the isomorphism above, it is a rational section of $\omega_{\mathcal Y_{\ell'}/R_{\ell'}}$ with divisor $p^e q_{\ell}V_{\ell'}-NV_{\ell'}$.
Taking divisors gives \eqref{eq:local-canonical-change}.
Since $p^e q_{\ell}<p^e r_{\ell}=r_{\ell'}$, the right-hand side of \eqref{eq:local-canonical-change} is greater than $-r_{\ell'}$.
The integer $n$ is therefore nonnegative.
The Jacobian differential extends unchanged through constant extension because the Jacobian has good reduction.
Equation \eqref{eq:hodge} then gives \eqref{eq:local-hodge-change}.
\end{proof}

\begin{proposition}\label{prop:local-normalization-stabilization}
Let $\mathcal Y/R$ be as in Lemma~\ref{lem:local-normalization}.
% the imperfection degree of $k$ is arbitrary.
For a field extension $\ell/k$, let $\nu_{\ell}\colon \mathcal Y_{\ell} \to \mathcal Y\times_R\Spec(R\otimes_k\ell)$ be the normalization.
Then there exists a finite extension $\ell/k$ such that $\nu_{\ell}$ is finite, $\mathcal Y_{\ell}$ is geometrically regular over $\ell$, and $V_{\ell}$ is geometrically reduced over $\ell$.
\end{proposition}

\begin{proof}
Fix a closed point $P$ of degree $d$ of the generic curve $Y$ over $K$.
For a finite extension $\ell/k$, let $K_{\ell}$ be the fraction field of $R_{\ell}$, and let $P_{\ell}$ be the unique point of $Y_{K_{\ell}}$ over $P$; its degree $d_{\ell}$ over $K_{\ell}$ is at most $d$.
Let $H_{\ell}$ be the closure of $P_{\ell}$ in the regular proper model $\mathcal Y_{\ell}$.
Then $H_{\ell}$ is a horizontal Cartier divisor, finite and flat of degree $d_{\ell}$ over $R_{\ell}$.
Intersecting it with the closed fibre gives
\[
  d_{\ell} = (H_{\ell}\cdot(r_{\ell}V_{\ell}))_{\ell} = r_{\ell}(H_{\ell}\cdot V_{\ell})_{\ell}.
\]
The final intersection is a positive integer, so $r_{\ell}\leq d$.

Choose $\ell_1$ for which $r_{\ell_1}$ is maximal among all finite purely inseparable constant extensions.
For every finite extension $\ell'/\ell_1$, maximality and \eqref{eq:local-multiplicity-change} give $r_{\ell'} = r_{\ell_1}$; hence $e(\ell''/\ell') = 0$ whenever $\ell_1\subseteq\ell'\subseteq\ell''$.
Choose $\ell_2\supseteq \ell_1$ for which $b(\mathcal Y_{\ell_2})$ is minimal among such extensions, which is possible because \eqref{eq:local-hodge-change} makes these nonnegative integers nonincreasing.
Hence $n(\ell''/\ell') = 0$ whenever $\ell_2\subseteq\ell'\subseteq\ell''$.
Finally choose $\ell_3\supseteq \ell_2$ for which $q_{\ell_3}$ is minimal.
For every $\ell'\supseteq \ell_3$, equation \eqref{eq:local-canonical-change} now gives
\[
  N(\ell'/\ell_3)=q_{\ell_3}-q_{\ell'}=0.
\]
The scheme $(\mathcal Y_{\ell_3}) \times_{R_{\ell_3}} \Spec R_{\ell'}$ is $S_2$ and regular in codimension one since $N(\ell'/\ell_3) = 0$, thus it is normal by Serre's criterion.
It consequently equals its normalization $\mathcal Y_{\ell'}$ and is regular by Lemma~\ref{lem:local-normalization}\,(1).
Since $(\mathcal Y_{\ell_3})_{R_{\ell'}} = \mathcal Y_{\ell_3}\times_{\ell_3}\Spec\ell'$ and every finite extension of $\ell_3$ is purely inseparable, $\mathcal Y_{\ell_3}$ is geometrically regular over $\ell_3$.
Since $e(\ell'/\ell_3) = 0$ for every finite extension $\ell'/\ell_3$, the pullback of $V_{\ell_3}$ to $\mathcal Y_{\ell_3}\times_{\ell_3}\Spec\ell' = \mathcal Y_{\ell'}$ is the reduced divisor $V_{\ell'}$; hence $V_{\ell_3}$ is geometrically reduced over $\ell_3$ \cite[Tag~\stackstag{035X}]{Stacks}.
Finally, $\nu_{\ell_3}$ is finite because $R_{\ell_3}$ is excellent.
\end{proof}

\subsection{Proof of Theorem~\ref{thm:local-hodge-congruence} and Corollary~\ref{lem:congruence}}\label{subsec:local-conductor-torsion}
\begin{proof}
We may replace $R$ by the completion of its strict henselization, which is isomorphic to $k^{\mathrm s}[[\pi]]$.
Indeed, this base change preserves the minimal regular model, its closed fibre, and generators of $H^0(\mathcal Y,\omega_{\mathcal Y/R})$, hence $r$, $q$, and $b$: for the strict henselization because it is ind-\'etale over $R$, and for the completion by Lemma~\ref{lem:ab-properties}(3).
After renaming $k^{\mathrm s}$ as $k$, Lemma~\ref{lem:local-normalization} and Proposition~\ref{prop:local-normalization-stabilization} apply.

Put $R_\infty=R\otimes_k\bar k$.
It is the union of the discrete valuation rings $R_{\ell'}$, for finite extensions $\ell'/k$, with uniformizer $\pi$, so every nonzero element of $R_\infty$ is a unit times a power of $\pi$.
Thus $R_\infty$ is a Noetherian discrete valuation ring.
By Proposition~\ref{prop:local-normalization-stabilization}, there is a finite extension $\ell/k$ such that $\nu_{\ell}$ is finite, $\mathcal Y_{\ell}$ is geometrically regular over $\ell$, and $V_{\ell}$ is geometrically reduced over $\ell$.

Fix such an $\ell$.
Then $\mathcal Z = \mathcal Y_{\ell}\times_{\ell}\Spec\bar k = \mathcal Y_{\ell}\times_{R_{\ell}}\Spec R_\infty$ is finite and birational over $\mathcal Y_{R_\infty}$; in particular, $\mathcal Z$ is Noetherian, since $\mathcal Y_{R_\infty}$ is of finite type over the Noetherian ring $R_\infty$.
Since $\mathcal Y_{\ell}$ is geometrically regular over $\ell$ and $\mathcal Z$ is Noetherian, $\mathcal Z$ is regular \cite[Proposition~5.13.7]{EGA}.
Being regular, finite and birational over $\mathcal Y_{R_\infty}$, the scheme $\mathcal Z$ is the normalization of $\mathcal Y_{R_\infty}$.

Since $V_{\ell}$ is geometrically reduced over $\ell$, the closed fibre of $\mathcal Z$ is $mE$, where $m = r_{\ell}$ and $E = V_{\ell}\times_{\ell}\Spec\bar k$ is a smooth elliptic curve.
By \eqref{eq:local-multiplicity-change}, $m = p^e r$ with $e = e(\ell/k)$.
Let $NE$ denote the divisorial conductor of $\mathcal Z\to \mathcal Y_{R_\infty}$.
It is the base change of the divisorial conductor $N(\ell/k)V_{\ell}$ of $\nu_{\ell}$, so $N = N(\ell/k)\geq0$ and $(p-1)\mid N$ by Lemma~\ref{lem:local-normalization}\,(2).
Let $\bar q$ be the integer in \eqref{eq:local-canonical-order} for $\mathcal Z$, and put $\bar a=a(\mathcal Z) = b(\mathcal Z)$, where the last equality holds by Lemma~\ref{lem:ab-properties}(4) since $\bar k$ is perfect.
The relations of Lemma~\ref{lem:local-normalization}\,(2) give
\begin{equation}\label{eq:local-conductor-comparison}
  nm=N+\bar q-p^e q,
  \qquad b(\mathcal Y)=n+\bar a.
\end{equation}

From now on, work over the completion $\bar k[[\pi]]$ of $R_\infty$.
The closed fibre $mE$ contains no $(-1)$-curve, so $\mathcal Z$ is the minimal regular model of its generic fibre; by Lemma~\ref{lem:ab-properties}(3) and flat base change, its base change to $\bar k[[\pi]]$ is again the minimal regular model, with the same closed fibre $mE$ and the same integers $\bar q$ and $\bar a$.
We replace $\mathcal Z$ by this base change.
By Proposition~\ref{prop:raynaud-chi} for $\mathcal Z$, whose integer $\chi$ is $m\bar a+\bar q$, we have $m\bar a+\bar q\equiv m-1\pmod{p-1}$.
Combining this with \eqref{eq:local-conductor-comparison} gives
\[
  m b(\mathcal Y)+p^e q
  =N+\bar q+m\bar a
  \equiv m-1\pmod{p-1}.
\]
Since $m=p^e r$ and $p^e\equiv1\pmod{p-1}$, this is \eqref{eq:local-hodge-congruence}.

Suppose that $r=1$.
Then $q=0$, so $(p-1)\mid b(\mathcal Y)$.
If $b(\mathcal Y)=0$, equation \eqref{eq:local-conductor-comparison} forces $n=\bar a=N=\bar q=0$.
The conductor vanishes, so $\mathcal Z=\mathcal Y\otimes_R\bar k[[\pi]]$.
Moreover, the integer $\chi = m\bar a+\bar q$ of $\mathcal Z$ vanishes; since every summand in \eqref{eq:raynaud-chi} for $\mathcal Z$ is nonnegative and the first one is $m-1$, we get $m=1$.
The closed fibre is smooth, and the generic fibre is smooth, so $\mathcal Z$ is smooth over $\bar k[[\pi]]$, and hence $\mathcal Y$ is smooth over $R$ by faithfully flat descent.
Conversely, if $\mathcal Y$ is smooth, then $b(\mathcal Y)=0$ by Lemma~\ref{lem:ab-properties}(2).
\end{proof}
